% ch8_a 第8章 §8.1–§8.2 (书页 335–344 / p350–p359)
% 第8章 TOPOLOGICAL AND QUANTUM ORDER—BEYOND LANDAU'S THEORIES

\chapter{拓扑序与量子序——超越朗道理论}

在第 3、4、5 章中，我们详细讨论了几个相互作用的玻色子/费米子系统。这些简单的模型例示了朗道的对称性破缺理论（加上 RG 理论）和朗道的费米液体理论（加上微扰论）。这两套朗道理论解释了许多凝聚态系统的行为，构成了传统多体理论的基础。这三章的讨论仅触及了朗道理论及其丰富应用的表面。想进一步学习朗道理论（以及 RG 理论和微扰论）的读者，会发现 Abrikosov et al. (1975)、Ma (1976)、Mahan (1990)、Negele and Orland (1998) 以及 Chaikin and Lubensky (2000) 的书颇有帮助。

朗道的理论非常成功，以至于在很长一段时间里，人们找不到任何不能用朗道理论描述的凝聚态系统。五十年来，朗道理论统治着多体物理，并实质上定义了多体物理的范式。这么多年之后，一种普遍的信念业已形成：我们已经弄清了所有重要的概念，理解了物质一切形态的基本性质。多体理论已经走到尽头，成为一门多少算是完备的理论。剩下要做的，只是把朗道理论（加上重整化群图像）应用于各种不同的系统。

从这个视角来看，我们能够理解 Tsui et al. (1982) 所发现的分数量子霍尔（FQH）效应的重要性。FQH 效应开辟了凝聚态物理的新篇章。正如我们在上一章看到的，FQH 液体不能用费米液体理论描述。不同的 FQH 态具有相同的对称性，因而不能用朗道的对称性破缺理论描述。因此，FQH 态完全超出了两套朗道理论。FQH 液体的存在表明，在朗道理论的范式之外还存在一个新世界。近来的研究表明，这个新范式比朗道理论的范式丰富得多。本书第 7 至 10 章将致力于讨论朗道理论之外的新范式。

为了窥见凝聚态物理新范式之一斑，除 FQH 态之外，我们还将研究量子转子系统、硬芯玻色子系统和量子自旋系统。这些系统都是强关联的多体系统。通常，这些系统中的强关联会导致长程序和对称性破缺。然而，我们已经以玻色超流体和费米子 SDW 态为例讨论过长程序和对称性破缺。因此，为了研究朗道理论之外的世界，我们把注意力集中于那些不能用长程序和对称性破缺描述的量子液体态（例如 FQH 态）。

这些量子液体态代表着物质的新状态，其中包含一种全新的序——拓扑序/量子序。我将证明，这种新序可能对我们理解量子相与量子相变以及量子相中的无能隙激发产生深远的影响。特别地，拓扑序/量子序或许能为自然界中的光和电子（以及其他规范玻色子和费米子）提供一个起源。

在本章中，我们将对拓扑序/量子序作一般性的讨论，以描绘一幅更大的图景。我们将以 FQH 态和费米液体态为例，讨论拓扑序/量子序中的一些基本问题。拓扑序/量子序中的问题与议题和传统多体物理中的大不相同。在深入任何细致的计算之前，理解这些问题究竟是什么是非常重要的。

\section{物质的状态与序的概念}

\begin{keypoints}
\item 物质可以有许多不同的状态（或不同的相）。序的概念是为了表征物质不同状态中不同的内部结构而引入的。
\item 我们曾以为，不同的序由其不同的对称性来表征。
\end{keypoints}

在足够高的温度下，一切物质都以气体的形式存在。气体是最简单的态之一。气体中一个原子的运动几乎不依赖于其他分子的位置和运动。因此，气体是弱关联系统，不含任何内部结构。然而，随着温度的降低，原子的运动变得越来越相互关联。最终，原子形成非常规则的图案，晶体序随之发展起来。在晶体中，单个原子几乎无法自行移动。晶体中的激发总是对应于许多原子的集体运动（这些集体运动称为声子）。晶体是强关联态的一个例子。

随着 1900 年前后低温技术的发展，物理学家发现了许多新的物质态（例如超导体和超流体）。这些不同的态具有不同的内部结构，称为不同种类的序。序的精确定义涉及相变。如果我们能够（通过平滑地改变哈密顿量）把一个多体系统的两个态平滑地相互变换而不遇到相变（即不遇到自由能的奇点），那么这两个态就具有相同的序。如果不存在不经相变就把一个态变为另一个态的途径，那么这两个态就具有不同的序。我们注意到，我们的序的定义是一个等价类的定义。能够不经相变而相互连通的两个态被定义为等价。以这种方式定义的等价类称为普适类。具有不同序的两个态也可以说是属于不同普适类的两个态。按照我们的定义，水和冰具有不同的序，而水和水蒸气具有相同的序（见图 8.1）。

%FIG{8.1}: {水的相图。}

在发现了如此多种不同的序之后，需要一个普遍理论以更深入地理解物质的态。特别地，我们想要理解，究竟是什么使两种序真正不同，以至于我们无法在不遇到相变的情况下把一种序变成另一种序。发展关于序以及相关的相与相变的普遍理论的关键一步，是认识到序与对称性相联系（更确切地说，与对称性的破缺相联系）。我们发现，当两个态具有不同的对称性时，我们无法在不遇到自由能奇点的情况下（即不经相变地）把一个态变为另一个态。基于序与对称性之间的这一关系，朗道发展了关于序以及不同序之间的相变的普遍理论（Ginzburg and Landau, 1950; Landau and Lifschitz, 1958）。朗道的理论非常成功。利用朗道理论和相关的对称性群论，我们可以对三维空间中可能存在的全部 230 种不同的晶体进行分类。通过确定对称性在连续相变处如何改变，我们可以得到相变的临界性质。对称性破缺还提供了许多无能隙激发的起源，例如声子、自旋波等，它们决定了许多系统的低能性质（Nambu, 1960; Goldstone, 1961）。这些激发的许多性质（包括其无能隙性）都直接由对称性决定。通过引入与对称性相联系的序参量，Ginzburg 和朗道发展了 Ginzburg--Landau 理论，它成为关于相与相变的标准理论。由于朗道的对称性破缺理论对我们理解物质具有如此广泛而根本的影响，它成了凝聚态理论的基石。朗道理论所描绘的图景令人如此满意，以至于人们开始有这样的感觉：至少在原则上，我们理解物质所能具有的一切种类的序。

\section{分数量子霍尔态中的拓扑序}

\begin{keypoints}
\item FQH 态开辟了凝聚态物理的新篇章，因为它是实验学家发现的第二个无法用对称性破缺和（局域）序参量来表征的态。
\item FQH 态包含一种新型的序——拓扑序。
\end{keypoints}

然而，自然界从不停止给我们带来惊喜。随着半导体技术的进步，物理学家学会了如何把电子限制在两种不同半导体之间的界面上，从而制成了二维电子气（2DEG）。1982 年，Tsui et al. (1982) 把二维电子气置于强磁场中，发现了一种新的物质态——FQH 液体（Laughlin, 1983）。由于温度低且电子之间的相互作用强，FQH 态是一个强关联态。然而，这样一个强关联态并不像人们原先预期的那样是晶体。事实证明，电子由于其质量非常之小而产生的强烈量子涨落阻止了晶体的形成。因此，FQH 态是一种量子液体。（晶体可以通过两种方式熔化：一是像我们升温时那样通过热涨落，得到普通的液体；二是像我们减小粒子质量时那样通过量子涨落，得到量子液体。）

正如我们在上一章看到的，量子霍尔液体具有许多令人惊叹的性质。量子霍尔液体比固体（晶体）更为“刚硬”，其含义是量子霍尔液体不可压缩。因此，量子霍尔液体具有固定且明确界定的密度。当我们用由下式定义的填充因子来度量电子密度时，
\begin{equation*}
\nu = \frac{\text{电子的密度}}{\text{磁通量子的密度}},
\end{equation*}
我们发现，所有已发现的量子霍尔态的密度都使填充因子恰好等于某些有理数，例如 $\nu = 1,\allowbreak 1/3,\allowbreak 2/3,\allowbreak 2/5,\allowbreak \ldots$。知道了 FQH 液体只存在于某些神奇的填充因子处，人们不禁猜测，FQH 液体应当具有某种内部的序或“图案”。不同的神奇填充因子应当来自这些不同的内部“图案”。然而，内部“图案”的假说看来有一个困难——FQH 态是液体，而液体怎么能有任何内部“图案”呢？

%FIG{8.2}: {圆周上的粒子波具有量子化的波长。}

为了对 FQH 态中的内部序获得一些直观理解，让我们尝试设想 FQH 态中电子的量子运动。我们知道，根据量子物理，粒子也表现出波的行为。我们先考虑一个在圆周上以动量 $p$ 运动的粒子。这样一个粒子对应于波长为 $\lambda = h/p$ 的波，其中 $h$ 是普朗克常数。只有能够纳入圆周的波才是允许的（即圆周必须包含整数个波长）（见图 8.2）。于是，由于量子物理，粒子在圆周上的运动受到很强的限制（即量子化），只有某些分立的动量值才是允许的。这个量子化条件可以用更具图像感的方式来看待。我们可以说，粒子沿着圆周以步伐跳舞，步长由波长给出。量子化条件要求粒子总是用整数步走完圆周一圈。

现在让我们考虑磁场中的一个电子。在磁场的作用下，电子总是沿圆周运动（称为回旋运动）。在量子物理中，由于粒子的波动性质，只有某些分立的回旋运动是允许的。量子化条件要求：允许的回旋运动的圆形轨道包含整数个波长。我们可以说，电子总是用整数步走完圆周一圈。如果电子绕圆周走了 $n$ 步，我们就说电子处在第 $n$ 个朗道能级。第一朗道能级中的电子能量最低，低温下电子将停留在第一朗道能级。

当我们有许多电子构成二维电子气时，电子不仅在第一朗道能级中作各自的回旋运动，而且还会彼此绕行并交换位置。这些附加的运动同样服从量子化条件。例如，一个电子必须用整数步绕另一个电子一周。由于电子是费米子，交换两个电子会给波函数引入一个负号。而交换两个电子等价于让一个电子绕另一个电子走半程。因此，一个电子必须用半整数步绕另一个电子走半程。（半整数步给电子波函数引入一个负号。）换言之，一个电子必须用奇数步绕另一个电子一周。FQH 态中的电子不仅以满足量子化条件的方式运动，而且由于电子之间强烈的库仑排斥和费米统计，它们还尽可能地彼此远离。这意味着，只要可能，一个电子就试图用更多的步数绕另一个电子一周。

现在我们看到，尽管不存在晶体序，FQH 态中电子的量子运动却是高度有组织的。FQH 态中的所有电子集体起舞，遵循严格的跳舞规则。

\begin{enumerate}
\item 所有电子都在第一朗道能级中作各自的回旋运动，即它们用一步绕圆周一圈。
\item 一个电子总是用奇数步绕另一个电子一周。
\item 电子试图彼此远离，即它们试图用尽可能多的步数绕另一个电子一周。
\end{enumerate}

如果每个电子都遵循这些严格的跳舞规则，那么只允许唯一的一种整体跳舞图案。这样的跳舞图案描述了 FQH 态中的内部量子运动。正是这个整体跳舞图案对应于 FQH 态中的内部序。不同的 FQH 态由其不同的跳舞图案相互区分。

对上文概述的电子量子运动，更精确的数学描述由著名的 Laughlin 波函数给出（Laughlin, 1983）
\begin{equation*}
\Psi_m = \Bigl[\prod (z_i - z_j)^m\Bigr]\, \mathrm{e}^{-\frac{1}{4l_B^2}\sum |z_i|^2}
\end{equation*}
其中 $m$ 是奇整数，$z_j = x_j + i y_j$ 是第 $j$ 个电子的坐标。这样的波函数描述填充因子为 $\nu = 1/m$ 的 FQH 态。我们看到，当 $z_i \rightarrow z_j$ 时波函数变为零，因此电子不喜欢彼此靠得太近。此外，当我们把一个电子绕另一个电子移动一周时，波函数的相位改变 $2\pi m$。因此，在 Laughlin 态中，一个电子总是用 $m$ 步绕另一个电子一周。

我们想强调，FQH 液体的内部序（即跳舞图案）与其他关联系统（如晶体、超流体等）中的内部序非常不同。后一类系统中的内部序可以用与破缺对称性相联系的序参量来描述，因而有序态可以用 Ginzburg--Landau 有效理论来描述。FQH 液体中的内部序是一种新型的序，无法用与破缺对称性相联系的长程序来描述。\footnote{尽管有人建议用“非对角长程序”（Girvin and MacDonald, 1987）来表征 Laughlin 态的内部结构，但具有长程序的那个算符本身并不是局域算符。对于局域算符而言，既不存在长程序，也不存在对称性破缺的 Laughlin 态。}1989 年，人们引入了“拓扑序”的概念来描述 FQH 液体中这种新型的序（Wen, 1990, 1995）。

我们想指出，拓扑序是零温下任何具有有限能隙的态的普遍性质。非平庸的拓扑序不仅出现在 FQH 液体中，也出现在零温下的自旋液体中。事实上，拓扑序的概念正是在手征自旋液体（Kalmeyer and Laughlin, 1987; Khveshchenko and Wiegmann, 1989; Wen et al., 1989）的研究中首次引入的（Wen, 1990）。除手征自旋液体外，在任意子超流体（Chen et al., 1989; Fetter et al., 1989; Wen and Zee, 1991）以及自旋系统的短程共振价键态（Kivelson et al., 1987; Rokhsar and Kivelson, 1988; Read and Chakraborty, 1989; Read and Sachdev, 1991; Wen, 1991a）中也发现了非平庸的拓扑序。FQH 液体甚至还不是实验上观察到的第一个具有非平庸拓扑序的态。这份荣誉属于 1911 年发现的超导态（Onnes, 1911; Bardeen et al., 1957）。与通常的观点相反，具有动力学电磁相互作用的超导态不能用破缺的对称性来表征。它既没有长程序，也没有局域序参量。超导态包含非平庸的拓扑序，它与超流态有着根本的不同（Coleman and Weinberg, 1973; Halperin et al., 1974; Fradkin and Shenker, 1979）。

把 FQH 液体与晶体作比较是富有启发性的。FQH 液体与晶体的相似之处在于，二者都包含丰富的内部图案（即内部序）。主要区别在于，晶体中的图案是静态的，与原子的位置相关；而 QH 液体中的图案是“动态的”，与电子彼此“跳舞”的方式相关。然而，关于晶体序可以提出的许多同样的问题，对拓扑序也应当提出并加以解决。我们知道，晶体序可以用对称性来表征和分类。于是，一个重要的问题是：我们如何表征和分类拓扑序？我们还知道，晶体序可以用 X 射线衍射来测量。第二个重要的问题是：我们如何在实验上测量拓扑序？

下面我们将更详细地讨论 FQH 态中的拓扑序。事实证明，FQH 态是相当典型的拓扑有序态。许多其他拓扑有序态与 FQH 态有许多相似的性质。

\subsection{拓扑序的表征}

\begin{keypoints}
\item 物理学中的任何新概念都必须借助实验可测的量来引入（或定义）。定义一个物理量或一个概念，就是设计一个实验。
\item 拓扑序的概念（部分地）由基态简并度定义，该简并度对能够破坏所有对称性的\emph{任何}微扰都是稳健的。
\end{keypoints}

在上文中，拓扑序的概念（跳舞图案）是通过基态波函数引入的。这并不十分正确，因为基态波函数不是普适的。要确立一个新概念（例如拓扑序），就需要找到拓扑序的物理表征或测量方法。换言之，需要找到普适的量子数，它们对任何微扰（诸如相互作用、有效质量等的改变）都稳健，但对不同类别的 FQH 液体可以取不同的值。这类量子数的存在蕴含着拓扑序的存在。

显示 FQH 液体中存在拓扑序的一种途径，是研究它们的基态简并度（在热力学极限下）。FQH 液体有一个非常特殊的性质：它们的基态简并度依赖于空间的拓扑（Haldane, 1983; Haldane and Rezayi, 1985）。例如，$\nu = 1/q$ 的 Laughlin 态在亏格为 $g$ 的黎曼面上具有 $q^g$ 个简并基态。FQH 液体中的基态简并并\emph{不}是哈密顿量的对称性的结果。基态简并度对任意微扰都稳健（甚至是破坏哈密顿量中所有对称性的杂质）（Wen and Niu, 1990）。基态简并度的稳健性表明，导致基态简并的内部结构是普适而稳健的，从而证明了普适内部结构——拓扑序——的存在。

为了理解 FQH 基态的拓扑简并，我们考虑环面上的 $\nu = 1/m$ Laughlin 态。我们将用两种方法来计算基态简并度。第一种方法考虑如下隧穿过程：我们先产生一个准粒子--准空穴对；然后让准粒子绕环面一整周；最后湮灭准粒子--准空穴对，回到基态。这样一个隧穿过程产生一个把基态映射到基态的算符。若准粒子沿 $x$ 方向绕环面一周，该算符记为 $U_x$；若沿 $y$ 方向绕环面一周，则记为 $U_y$（见图 8.3(a,b)）。于是，沿 $x$、$y$、$-x$、$-y$ 四个方向的四次隧穿生成 $U_y^{-1}U_x^{-1}U_yU_x$（见图 8.3(c)）。我们注意到，上述四次隧穿的路径可以变形为两个相互链接的回路（见图 8.3(d)）。这两个相互链接的回路对应于让一个准粒子绕另一个准粒子一周，它产生一个相位 $\mathrm{e}^{2\mathrm{i}\theta}$，其中 $\theta$ 是准粒子的统计角。对 $1/m$ Laughlin 态，我们有 $\theta = \pi/m$。因此我们得到
\begin{equation*}
U_y^{-1}U_x^{-1}U_yU_x = \mathrm{e}^{\mathrm{i}2\pi/m}
\end{equation*}

%FIG{8.3}: {(a) 沿 $x$ 方向的隧穿产生 $U_x$。(b) 沿 $y$ 方向的隧穿产生 $U_y$。(c) 沿 $x$、$y$、$-x$ 和 $-y$ 方向的四次隧穿产生 $U_y^{-1}U_x^{-1}U_yU_x$。(d) 上述四次隧穿可以变形为两个相互链接的回路。}

由于 $U_{x,y}$ 作用在基态之内，基态构成上述代数的表示。该代数只有一个 $m$ 维不可约表示。因此，$1/m$ Laughlin 态具有 $m \times$ 整数个简并基态。这一途径使我们能够直接看到准粒子统计与基态简并度之间的联系。

计算基态简并度的第二种方法是利用有效理论 (7.3.11)。简并的基态来自以下的集体涨落：
\begin{equation*}
a_i(t, x, y) = \theta_i(t)/L, \qquad i = x, y \tag{8.2.1}
\end{equation*}
其中 $L$ 是环面在 $x$ 和 $y$ 方向上的尺寸。所有其他涨落都产生非零的“磁”场 $b = f_{xy}$，并具有有限的能隙，这一点从经典运动方程可以看出。把式 (8.2.1) 代入式 (7.3.11)，可以得到描述式 (8.2.1) 中集体激发动力学的下列拉格朗日量：
\begin{equation*}
L = - \frac{m}{4\pi} (\dot{\theta}_x \theta_y - \dot{\theta}_y \theta_x) + \frac{1}{2g_1} \dot{\theta}_i^2. \tag{8.2.2}
\end{equation*}

由于 $a_\mu$ 的荷以整数量子化，作用在准粒子场 $\psi_q \to U\psi_q$ 上的规范变换 $U(x,y)$ 必须是环面上的周期函数。因此，规范变换必具有形式 $U(x,y) = \exp\bigl(2\pi\mathrm{i}\bigl(\frac{nx}{L} + \frac{my}{L}\bigr)\bigr)$，其中 $n$ 和 $m$ 是整数。由于 $\psi_q$ 的 $a_\mu$ 荷为 1，这样的规范变换按下式把规范场 $a_i$ 变为 $a_i' = a_i - \mathrm{i}U^{-1}\partial_i U$：
\begin{equation*}
(a_x', a_y') = \Bigl(a_x + \frac{2\pi n}{L},\ a_y + \frac{2\pi m}{L}\Bigr) \tag{8.2.3}
\end{equation*}

式 (8.2.3) 蕴含，$(\theta_x, \theta_y)$ 与 $(\theta_x + 2\pi n, \theta_y + 2\pi m)$ 规范等价，应当认同为同一位形。规范不等价的位形由环面 $0 < \theta_i < 2\pi$ 上的一个点给出。于是，拉格朗日量 (8.2.2) 描述一个单位电荷的粒子在以 $(\theta_x, \theta_y)$ 为参数的环面上运动。

式 (8.2.2) 的第一项表明，环面上存在均匀的“磁”场 $B = m/2\pi$。穿过环面的总通量等于 $2\pi \times m$。式 (8.2.2) 的哈密顿量为
\begin{equation*}
H = \frac{g_1}{2} \Bigl[- (\partial_{\theta_x} - \mathrm{i} A_{\theta_x})^2 - (\partial_{\theta_y} - \mathrm{i} A_{\theta_y})^2 \Bigr]. \tag{8.2.4}
\end{equation*}

式 (8.2.4) 的能量本征态形成朗道能级。朗道能级之间的能隙为 $g_1$ 的量级，与系统的尺寸无关。第一朗道能级中的态的数目等于穿过环面的磁通量子的数目，在我们的情形是 $m$。因此，$1/m$ Laughlin 态的基态简并度为 $m$。

为了理解基态简并度的稳健性，让我们在电子哈密顿量上加一个任意微扰。这样的微扰将引起有效拉格朗日量的改变 $\delta\mathcal{L}(a_\mu)$。这里的关键是，$\delta\mathcal{L}(a_\mu)$ 只通过场强依赖 $a_\mu$。换言之，$\delta\mathcal{L}$ 是 $e$ 和 $b$ 的函数。由于式 (8.2.1) 中的集体涨落在局域上是纯规范，$e$ 和 $b$ 从而 $\delta\mathcal{L}$ 都不依赖 $\theta_i$。$\delta\mathcal{L}(a_\mu)$ 不能在 $\theta_i$ 的有效理论中生成任何势项 $V(\theta_i)$；它只能生成仅依赖 $\dot{\theta}_i$ 的项。这样的修正重整化 $g_1$ 的数值，不能解除简并。

更一般的 FQH 态由式 (7.3.18) 描述。当 $K$ 为对角矩阵时，上述结果意味着基态简并度为 $\det(K)$。事实证明，对于一般的 $K$，基态简并度同样由 $\det(K)$ 给出。在亏格为 $g$ 的黎曼面上，基态简并度变为 $(\det(K))^g$。

我们看到，在紧致空间上，FQH 液体的低能物理非常独特。低能激发只有有限多个（即那些简并基态），然而低能动力学却是非平庸的，因为基态简并度依赖于空间的拓扑。这类只依赖于空间拓扑的特殊低能动力学由所谓的拓扑场论描述；拓扑场论在高能物理界曾被深入研究（Witten, 1989; Elitzur et al., 1989; Fr\"{o}hlich and King, 1989）。拓扑理论是 FQH 液体的有效理论，正如 Ginzburg--Landau 理论是超流体（或其他对称性破缺相）的有效理论。

基态简并度对空间拓扑的依赖表明，尽管所有局域物理算符都不具有长程关联，FQH 液体中仍存在某种长程序（即上面提到的整体跳舞图案）。在某种意义上，我们可以说 FQH 液体包含着隐藏的长程序。
